\section{Filtered quotient naturality}
\label{sec:quotient}

The transfer theorem determines the dimension of every endpoint kernel. Exact filling determines the maps between those kernels at the homology level. This separation is useful: spectral concentration proves that there are no extra classes, while the quotient model computes persistence without comparing harmonic representatives.

\begin{lemma}[Independent filling]
\label{lem:independent-filling}
For every term set \(A\),
\begin{equation}
\label{eq:boundary-intersection}
B_d(X_A)\cap V=W_A
=\sum_{j\in A}\ran\Pi_j.
\end{equation}
\end{lemma}

\begin{proof}
Every vector in \(\ran\Pi_j\) has a filling in gadget \(j\), so \(W_A\subseteq B_d(X_A)\cap V\). Conversely, write a degree-\((d+1)\) chain as a sum of its gadget components. If its total boundary lies in \(V\), then its boundary has zero coordinate on every private \(d\)-simplex. Private supports of distinct gadgets are disjoint, so this vanishing holds separately for each gadget component. The remaining register boundary of component \(j\) lies in \(B_d(Y_j)\cap V=\ran\Pi_j\). Summing over \(j\) proves the reverse inclusion.
\end{proof}

No relative-acyclicity hypothesis appears in \cref{lem:independent-filling}. Extra geometric homology outside \(V\) is excluded later by the all-chain concentration theorem, not by the boundary decomposition alone.

\begin{theorem}[Natural quotient realization]
\label{thm:quotient}
Under the hypotheses and parameter choice of \cref{thm:transfer}, the map
\begin{equation}
\label{eq:qA}
q_A:V/W_A\longrightarrow H_d(X_A),
\qquad
q_A(v+W_A)=[v],
\end{equation}
is an isomorphism. For \(A\subseteq B\), these isomorphisms make the diagram
\[
\begin{CD}
V/W_A @>{\rho_{AB}}>> V/W_B\\
@V{q_A}V{\cong}V @VV{\cong}V{q_B}\\
H_d(X_A) @>{(\iota_{AB})_*}>> H_d(X_B)
\end{CD}
\]
commute, where \(\rho_{AB}\) is the natural quotient map. Consequently,
\begin{equation}
\label{eq:persistent-rank}
\rank(\iota_{AB})_*
=\dim(V/W_B)
=\dim K_B.
\end{equation}
\end{theorem}

\begin{proof}
\Cref{eq:boundary-intersection} makes \(q_A\) well-defined and injective. Its domain has dimension
\[
\dim(V/W_A)=\dim K_A.
\]
By Hodge theory and \cref{eq:kernel-gap-conclusion}, the codomain has dimension
\(
\beta_d(X_A)=\dim\ker\Delta_A=\dim K_A
\).
Thus \(q_A\) is surjective as well.

If \(A\subseteq B\), both paths send \(v+W_A\) to the class of the same register cycle \(v\) in \(X_B\). The diagram therefore commutes. Since \(W_A\subseteq W_B\), \(\rho_{AB}\) is onto, and its rank equals the dimension of \(V/W_B\). Finally, \(W_B=\ran H_B=K_B^\perp\) inside \(V\), which proves~\cref{eq:persistent-rank}.
\end{proof}

For a chain \(A_0\subseteq A_1\subseteq\cdots\subseteq A_s\), every arrow in degree \(d\) is an epimorphism. The constructed barcode therefore has deaths but no new births in this degree. More generally, for any \(i\le j\),
\[
\rank\bigl[H_d(X_{A_i})\to H_d(X_{A_j})\bigr]
=\dim K_{A_j}.
\]
The result concerns the algebraic homology maps. Harmonic subspaces at different filtration levels need not be related by a literal common isometry; their Hodge representatives can rotate while the quotient maps remain exact.

Taking \(A_{\rm in}\) to be the clock, propagation, and clean-input terms of a history Hamiltonian, and taking \(A_{\rm out}\) to add the final rejection term, gives
\begin{equation}
\label{eq:history-ratio}
q_d(\Xin,\Xout)
=\frac{\dim K_{\rm out}}{\dim K_{\rm in}}.
\end{equation}
This identity is the bridge from exact logical perfect-space fractions to true normalized persistence.
